% TOP_PROBLEM: {"id": 30007510, "title": "Kervaire detection at the prime three", "category_id": 7, "category": {"id": 7, "name": "topology", "display_name": "Topology", "description": "Properties preserved under continuous deformations.", "slug": "topology", "order_index": 7, "created_at": "2026-07-31T15:26:25.671Z"}, "status": "open"}
% FIRST_POSTED: 2026-10-06
% ATTEMPT_STATUS: solved
% OPENAI_CLAIM: {"schema_version":1,"source_repo":"https://github.com/openai/math","source_commit":"adc7f1241b42e322a6451854ab7e4b4c146bf78a","release_date":"2026-10-06","match":"full","resolution":"proved","independently_reviewed":false,"second_pass":"agreed","families":[{"family":"309","title":"The Kervaire invariant problem at the prime three","subject":"Topology","manuscripts":[{"title":"The Kervaire invariant problem at the prime three","pdf_path":"preprints/The-Kervaire-Invariant-Problem-at-the-Prime-Three-September-24-2026/paper.pdf","readme_path":"preprints/The-Kervaire-Invariant-Problem-at-the-Prime-Three-September-24-2026/README.md","date":"2026-09-24","theorem_ref":"preprints/The-Kervaire-Invariant-Problem-at-the-Prime-Three-September-24-2026/build/sections/introduction.tex#thm:classification"}],"lean":null,"reasoning_trace":null}]}
\documentclass[11pt]{article}
\usepackage[margin=25mm]{geometry}
\usepackage{fontspec}
\setmainfont{Latin Modern Roman}
\usepackage{amsmath,amssymb,mathtools}
\usepackage{unicode-math}
\setmathfont{Latin Modern Math}
\usepackage{enumitem}
\usepackage{xurl,hyperref,bookmark}
\hypersetup{colorlinks=true,urlcolor=blue}
\setcounter{secnumdepth}{0}
\setlength{\parindent}{0pt}
\setlength{\parskip}{5pt}
\setlength{\emergencystretch}{3em}
\begin{document}
\section*{Kervaire detection at the prime three}

\subsection{OpenAI's claim}
The standard mod-three Adams Kervaire classes are claimed to survive exactly at indices zero, two, and three, each with an order-three representative; the new index-three element lies in stem 322. The cited manuscript argues for that result; its relation to this problem is specified below.

\subsection{Outline of the argument}
Manuscript 1 (family 309, 2026-09-24): A height-six Morava detector with C\_9 action imposes finite Euler deadlines and outgoing differentials on all higher candidate classes. Independently, a truncated extended cube of the known stem-106 Moore map carries the required primary operation. Finite obstruction corrections lift it back through the Moore sequence, preserving detection and producing an actual sphere element killed by three.

\subsection{Scope relative to this problem}
Manuscript 1 (family 309, 2026-09-24): The problem selects the mathematical statement summarized here, including its stated hypotheses and parameter range. The cited theorem gives a claimed proof of that assertion in full.

\subsection{Formal verification}
Family 309 has no Lean scope page in this snapshot; no formalization of these manuscript results is listed.

This site has not executed the Lean files or checked their dependencies. These descriptions report the scope of the linked declarations, rather than a new verification of the proofs.

\subsection{Status caveat}
OpenAI attributes these results to an internal model. Its release notice cautions that results lacking full formalization can contain mistakes. This site has not independently assessed the proofs; a solved label records the scope of the reported claim.

\subsection{Sources}
\begin{itemize}
\item OpenAI, \emph{The Kervaire invariant problem at the prime three}, 2026-09-24; citation key \texttt{OAI:The-Kervaire-Invariant-Problem-at-the-Prime-Three-September-24-2026}. \href{https://github.com/openai/math/blob/adc7f1241b42e322a6451854ab7e4b4c146bf78a/preprints/The-Kervaire-Invariant-Problem-at-the-Prime-Three-September-24-2026/paper.pdf}{PDF (GitHub)}; \href{https://github.com/openai/math/raw/adc7f1241b42e322a6451854ab7e4b4c146bf78a/preprints/The-Kervaire-Invariant-Problem-at-the-Prime-Three-September-24-2026/paper.pdf}{PDF (direct)}; \href{https://github.com/openai/math/blob/adc7f1241b42e322a6451854ab7e4b4c146bf78a/preprints/The-Kervaire-Invariant-Problem-at-the-Prime-Three-September-24-2026/README.md}{Manuscript page and citation}; \href{https://github.com/openai/math/blob/main/preprints/The-Kervaire-Invariant-Problem-at-the-Prime-Three-September-24-2026/paper.pdf}{Latest version}.
\item \href{https://github.com/openai/math/blob/adc7f1241b42e322a6451854ab7e4b4c146bf78a/CONTENTS.md}{OpenAI catalogue entry}.
\end{itemize}
\end{document}
